\documentclass[11pt]{article}

\usepackage[margin=1in]{geometry}
\usepackage{booktabs}
\usepackage{longtable}
\usepackage{array}
\usepackage{hyperref}
\usepackage{url}
\usepackage{amsmath}

\title{Piano Transcription Research: Experiment Log}
\author{Project research record}
\date{Updated 8 November 2026}

\begin{document}
\maketitle

\section{Purpose and status}

This log records the reproducibility work, code changes, fine-tuning runs,
held-out evaluation, inference ablations, and timing-error diagnostics for the
structured-attribute piano transcription experiments in this repository. The
main question is whether adding structured velocity and timing attributes to
the TransKun transcription model improves transcription quality while
preserving the pretrained note-recognition baseline.

\paragraph{Current status.}
The 10,000-step fine-tuning stage is the strongest run evaluated so far.
Reducing the structured-attribute contribution at inference did not improve
SMD scores. Matched-note timing diagnostics show small median offset bias but
substantial offset errors on some notes, particularly for longer reference
durations. A paired comparison on fixed MAESTRO validation windows found no
demonstrated improvement over the original checkpoint; the intervals remain
wide enough that this is inconclusive rather than proof of equivalence or
regression. A paired full-recording comparison is now being resumed from
saved per-piece results. It is an exploratory robustness analysis on the
same validation split, not independent confirmation; its partial results
must not be interpreted as an outcome. Do not select hyperparameters on SMD
going forward: this 50-piece set has already been used for iterative
comparisons. Reserve the untouched MAESTRO test set for one final evaluation
only after the method and decision rule are frozen.

\section{Evaluation conventions}

\begin{itemize}
  \item \textbf{SMD set:} 50 recordings under
        \texttt{datasets/SMD-piano\_v2}; all evaluations below use matching
        44.1 kHz stereo WAVs and MIDI references.
  \item \textbf{Ground-truth pedal extension:} disabled
        (\texttt{--noPedalExtension}).
  \item \textbf{Onset tolerance:} 50 ms
        (\texttt{--onsetTolerance 0.05}).
  \item \textbf{Reported SMD scores:} macro-averaged per-recording
        precision, recall, F1, and overlap from
        \texttt{transkun.computeMetrics}; the tables show F1.
  \item \textbf{Metric names:}
        \begin{itemize}
          \item \texttt{note}: onset and pitch.
          \item \texttt{note+offset}: onset, pitch, and note offsets.
          \item \texttt{note+velocity}: onset, pitch, and velocity.
          \item \texttt{note+velocity+offset}: all three.
        \end{itemize}
  \item \textbf{Timing error sign:} residuals are
        reference time minus estimated time. A negative onset or offset
        residual means the estimate is late; a negative duration residual
        means the estimated duration is longer than the reference.
\end{itemize}

Unless stated otherwise, the model uses
\texttt{transkun/pretrained/2.0.conf}, batch size 1, one training process,
two DataLoader workers, and the prepared 44.1 kHz MAESTRO copy. The two-worker
configuration succeeded after an earlier Windows paging-file failure with
the first pilot attempt.

The standard SMD score command was:

\begin{verbatim}
$py = ".\.venv\Scripts\python.exe"
$estimate = ".\checkpoint\structured-attributes-10000-smd\midi"
$truth = ".\datasets\SMD-piano_v2\midi"
$result = ".\checkpoint\structured-attributes-10000-smd\metrics.json"
& $py -m transkun.computeMetrics $estimate $truth `
  --noPedalExtension --onsetTolerance 0.05 --nProcess 1 `
  --outputJSON $result
\end{verbatim}

\section{Code and reproducibility work}

\begin{longtable}{p{0.19\linewidth}p{0.73\linewidth}}
\toprule
\textbf{Commit} & \textbf{Work recorded} \\
\midrule
\endhead
\texttt{03b23776} &
Fixed single-process MIDI evaluation argument handling, onset alignment
when no notes match, consistent onset tolerance, and regression tests. \\
\texttt{354b4c4d} &
Added structured velocity, timing, and presence factors to interval
transcription, configuration, documentation, and tests. \\
\texttt{dd035e7a} &
Added compatible pretrained-weight initialization for fine-tuning the
structured-attribute model. \\
\texttt{bfbb2953} &
Improved the MAESTRO fine-tuning workflow and terminal reporting; added
UTF-8 manifest handling, a finite optimizer-step limit, zero-worker
DataLoader support, documentation, and a MAESTRO manifest test. \\
\bottomrule
\end{longtable}

The original packaged checkpoint contains 184 state entries: 176 compatible
backbone/scorer tensors transfer to the new model, four structured-attribute
head tensors start from their neutral initialization, and eight old
velocity/refined-offset tensors are not reused. The new head must therefore
learn its attributes during fine-tuning; the original checkpoint's
velocity-inclusive SMD score is not directly comparable to an untrained new
head.

The MAESTRO source manifest contained 1,276 entries. The prepared training
copy resampled 233 48 kHz WAVs to 44.1 kHz and hard-linked 2,319 unchanged
audio/MIDI files; the source data was preserved. No-pedal-extension metadata
contains 962 training, 137 validation, and 177 test entries. A real sample
loaded as a 16-second, two-channel audio slice and produced gradients in the
structured-attribute head.

The validation metadata contains 69,946.2 seconds (19.43 hours) across 137
pieces and 639,425 positive-pitch note events. Piece duration has a median
of 359.5 seconds and a mean of 510.6 seconds; the 90th percentile is
1,113.7 seconds and the maximum is 2,457.8 seconds. Positive-pitch note
counts have a median of 2,907 per piece, a mean of 4,667, and a maximum of
23,469. Pedal events are excluded from this note count. These descriptive
values were computed directly from the prepared validation metadata and
motivate reporting macro-averaged per-piece scores.

The repository unit suite passed: \textbf{9 tests}. Coverage includes two
single-process evaluation cases, UTF-8 MAESTRO manifest parsing, two
pretrained-initialization cases, and four structured-attribute cases.

\section{Original-model SMD reference}

The original checkpoint-era model was evaluated on all 50 SMD recordings,
without ground-truth pedal extension and with 50 ms onset tolerance. Three
of four reproduced F1 values closely match the repository's published SMD
No Ext values; note+velocity differs by about 0.008 and needs reconciliation
before a formal comparison.

\begin{center}
\begin{tabular}{lrr}
\toprule
\textbf{Metric} & \textbf{Reproduced F1} & \textbf{README SMD No Ext F1} \\
\midrule
Note onset & 0.987627 & 0.9876 \\
Note onset + offset & 0.785080 & 0.7855 \\
Note + velocity & 0.953269 & 0.9455 \\
Note + velocity + offset & 0.758407 & 0.7588 \\
\bottomrule
\end{tabular}
\end{center}

The note+velocity difference is recorded rather than attributed to a
specific cause; evaluation settings, checkpoint version, or aggregation may
differ and should be checked before interpreting it.

\section{Fine-tuning experiments}

\subsection{1,000-step pilot}

The pilot fine-tuned from the packaged V2 checkpoint for 1,000 optimizer
steps. The first attempt with two DataLoader workers stopped before learning
because Windows ran out of paging-file capacity (\texttt{WinError 1455}).
It was restarted using \texttt{--dataLoaderWorkers 0}; validation then
completed. The final pilot checkpoint records step 1,000.

\begin{verbatim}
$py = ".\.venv\Scripts\python.exe"
$data = ".\datasets\maestro-v3.0.0-44100"
$train = "$data\metadata-no-pedal-extension\train.pickle"
$val = "$data\metadata-no-pedal-extension\val.pickle"
& $py -u -m transkun.train .\checkpoint\structured-attributes-pilot.pt `
  --nProcess 1 --datasetPath $data `
  --datasetMetaFile_train $train --datasetMetaFile_val $val `
  --modelConf .\transkun\pretrained\2.0.conf --nIter 1000 `
  --batchSize 1 --dataLoaderWorkers 0 `
  --pretrainedCheckpoint .\transkun\pretrained\2.0.pt
\end{verbatim}

\begin{center}
\begin{tabular}{lrrrr}
\toprule
\textbf{MAESTRO validation} & \textbf{Mean NLL} & \textbf{Precision} &
\textbf{Recall} & \textbf{F1} \\
\midrule
1k pilot & 111.2404 & 0.7851 & 0.7459 & 0.7650 \\
5k stage & 98.9169 & 0.8040 & 0.7620 & 0.7825 \\
10k stage & 97.5721 & 0.8054 & 0.7658 & 0.7851 \\
\bottomrule
\end{tabular}
\end{center}

Validation F1 here is the training validation's note+offset matching score.
It is not the velocity-inclusive SMD score. The values for the later stages
were recovered from their TensorBoard event files rather than inferred from
rounded F1 values in the experiment summary. From 1k to 5k, F1 increased by
0.0175 and mean NLL decreased by 12.3235; from 5k to 10k, F1 increased by
0.0026 and mean NLL decreased by 1.3448. These are sequential checkpoints,
not independent training replications.

\subsection{SMD comparison by training stage}

Each later stage was initialized from the preceding stage's checkpoint with
a fresh learning-rate scheduler. Thus the stages used 1,000, then 5,000,
then 10,000 additional optimizer updates (16,000 stage-wise updates in total
after the packaged initialization). The checkpoint step counter is local to
each stage.

\begin{verbatim}
$py = ".\.venv\Scripts\python.exe"
$data = ".\datasets\maestro-v3.0.0-44100"
$train = "$data\metadata-no-pedal-extension\train.pickle"
$val = "$data\metadata-no-pedal-extension\val.pickle"
& $py -u -m transkun.train .\checkpoint\structured-attributes-5000.pt `
  --nProcess 1 --datasetPath $data `
  --datasetMetaFile_train $train --datasetMetaFile_val $val `
  --modelConf .\transkun\pretrained\2.0.conf --nIter 5000 `
  --batchSize 1 --dataLoaderWorkers 2 `
  --pretrainedCheckpoint .\checkpoint\structured-attributes-pilot.pt

& $py -u -m transkun.train .\checkpoint\structured-attributes-10000.pt `
  --nProcess 1 --datasetPath $data `
  --datasetMetaFile_train $train --datasetMetaFile_val $val `
  --modelConf .\transkun\pretrained\2.0.conf --nIter 10000 `
  --batchSize 1 --dataLoaderWorkers 2 `
  --pretrainedCheckpoint .\checkpoint\structured-attributes-5000.pt
\end{verbatim}

\begin{center}
\begin{tabular}{lrrrr}
\toprule
\textbf{SMD F1} & \textbf{1k pilot} & \textbf{5k stage} &
\textbf{10k stage} & \textbf{Original model} \\
\midrule
Note onset & 0.983513 & 0.986316 & 0.986832 & 0.987627 \\
Note onset + offset & 0.752808 & 0.779505 & 0.779944 & 0.785080 \\
Note + velocity & 0.786829 & 0.951087 & 0.952227 & 0.953269 \\
Note + velocity + offset & 0.630164 & 0.751711 & 0.752384 & 0.758407 \\
\bottomrule
\end{tabular}
\end{center}

The large 1k-to-5k increase in velocity-inclusive scores is expected because
the structured-attribute head was newly initialized. The 5k-to-10k gains are
small. Validation F1 was 0.7825 after the 5k stage and 0.7851 after the 10k
stage.

\paragraph{Paired 5k versus 10k comparison.}
On per-recording note+offset F1, the 10k stage averaged +0.000439 over 5k;
24 pieces improved and 26 worsened. A piece-resampled bootstrap 95\%
interval was [-0.003672, +0.004996], consistent with little or no reliable
per-piece gain. The other mean paired F1 deltas were +0.000517 for note
onset, +0.001141 for note+velocity, and +0.000674 for
note+velocity+offset. Results use the same 50 SMD pieces and should be
treated as exploratory.

\section{Inference-only attribute-weight ablation}

Using the 10k checkpoint, the structured-attribute contribution to
semi-CRF interval scoring was varied without changing model weights. The
default weight is 1.0.

\begin{center}
\begin{tabular}{lrrrr}
\toprule
\textbf{Weight} & \textbf{Note} & \textbf{Note+offset} &
\textbf{Note+velocity} & \textbf{Note+velocity+offset} \\
\midrule
0.0  & 0.926784 & 0.715898 & 0.892112 & 0.689540 \\
0.5  & 0.984944 & 0.774892 & 0.950578 & 0.747845 \\
0.75 & 0.986230 & 0.778220 & 0.951585 & 0.750646 \\
1.0  & 0.986832 & 0.779944 & 0.952227 & 0.752384 \\
\bottomrule
\end{tabular}
\end{center}

Weight 1.0 scored higher overall than 0.5 or 0.75. For note+offset F1, the
paired mean differences against weight 1.0 were -0.001723 at weight 0.75
(bootstrap 95\% interval [-0.002335, -0.001077]) and -0.005052 at weight
0.5 ([-0.006206, -0.003947]). Weight 1.0 also beat weight 0.0 on all 50
pieces for each reported metric. This supports retaining the default weight;
it does not establish that 1.0 is optimal on unseen data.

\section{Offset timing diagnostic}

The saved JSON \texttt{deviations} values are not used for this diagnostic:
the repository computes those pairs with 0.8 s onset and minimum-offset
tolerances, which are much looser than the reported evaluation. Instead, the
diagnostic reloads predicted and reference MIDI and pairs notes using
\texttt{mir\_eval} rules: 50 ms onset tolerance; for onset+offset matching,
offset tolerance is the greater of 20\% of reference duration and 50 ms.
Ground-truth pedal extension is disabled. Results below are therefore
conditional on matched notes, not an error distribution over missed notes.

\subsection{Matched-note residuals}

\begin{center}
\begin{tabular}{llrr}
\toprule
\textbf{Residual, reference-estimate} & \textbf{Subset} &
\textbf{5k} & \textbf{10k} \\
\midrule
Onset median & Onset+pitch matched & 0.00 ms & 0.00 ms \\
Onset mean absolute error & Onset+pitch matched & 5.45 ms & 5.41 ms \\
Offset median & Onset+pitch matched & -3.65 ms & -2.08 ms \\
Offset median absolute error & Onset+pitch matched & 16.15 ms & 16.15 ms \\
Offset within 50 ms & Onset+pitch matched & 76.03\% & 76.12\% \\
Offset median & Strict onset+offset matched & -3.13 ms & -2.60 ms \\
Offset median absolute error & Strict onset+offset matched & 11.46 ms & 11.46 ms \\
Duration median & Strict onset+offset matched & -4.17 ms & -3.65 ms \\
\bottomrule
\end{tabular}
\end{center}

For onset+pitch matched notes, the 10k checkpoint had 142,991 pairs.
The offset-error mean was -1.45 ms, but mean absolute error was 70.55 ms;
the 5th and 95th percentiles were -140.1 ms and +156.8 ms. The small signed
median therefore hides a broad, approximately balanced error tail. Negative
offset residual means the estimated note ends later than the reference.

The strict onset+offset matching procedure found 114,090 matches. Separately,
applying the offset threshold to the onset-only matched pairs found 29,009
offset-threshold failures (20.29\%). These counts need not be exact
complements because maximum bipartite assignment can select different pairs
when offset constraints are enabled.

\subsection{Offset failures by reference duration}

Among onset+pitch matched notes, the offset failure rates varied by
reference duration:

\begin{center}
\begin{tabular}{lrrr}
\toprule
\textbf{Reference duration} & \textbf{Pairs} & \textbf{Offset fail rate} &
\textbf{Mean absolute offset error} \\
\midrule
Under 250 ms & 109,747 & 14.51\% & 32.22 ms \\
250--500 ms & 17,906 & 42.09\% & 117.61 ms \\
500 ms--1 s & 9,243 & 39.68\% & 204.86 ms \\
At least 1 s & 6,095 & 30.83\% & 418.89 ms \\
\bottomrule
\end{tabular}
\end{center}

Long notes are relatively few but have much larger absolute offset errors.
This points to duration-conditioned error analysis as the next research
direction rather than a global onset shift or reducing the attribute weight.
The duration bins and matching are descriptive; they do not prove a specific
mechanism.

\subsection{MAESTRO validation timing sample}

To check whether the SMD pattern generalizes beyond the repeatedly used
external comparison set, a separate diagnostic was run on the 137-piece
MAESTRO validation split using the 10k checkpoint. One uniformly sampled,
deterministic 16-second window was selected from each piece (Python random
seed 20261108). Notes were evaluated only in the center 12 seconds, leaving
two seconds of context on each side. Reference notes had to be strictly
contained in the window; ground-truth pedal extension was disabled.
Consequently this is a sampled-window diagnostic, not full-split validation
or a replacement for the training validation F1.

Among 16,500 reference notes and 15,857 estimated notes in the evaluated
regions, 15,776 notes matched on pitch and onset within 50 ms. Of those,
13,020 were found by the constrained onset+offset matching. Directly applying
the offset criterion to the onset-matched pairs gave 2,760 failures (17.49\%).
The small four-pair difference between this count and the constrained match
count is due to the matching algorithm being able to choose different
assignments when offset constraints are enabled.

\begin{center}
\small
\begin{tabular}{p{3.1cm}p{3.5cm}rrr}
\toprule
\textbf{Residual} & \textbf{Match subset} &
\textbf{Signed median} & \textbf{Median abs.} & \textbf{Mean abs.} \\
\midrule
Onset & Onset+pitch & -0.18 ms & 2.95 ms & 3.82 ms \\
Offset & Onset+pitch & +1.13 ms & 11.85 ms & 58.49 ms \\
Offset & Constrained onset+offset & +0.38 ms & 8.50 ms & 14.75 ms \\
Duration & Constrained onset+offset & +0.29 ms & 8.05 ms & 14.53 ms \\
\bottomrule
\end{tabular}
\normalsize
\end{center}

For onset+pitch matched notes, 80.13\% of offsets were within an absolute
50 ms, although the mean absolute error was much larger than the median. The
signed median remains close to zero; the principal issue is the error tail,
not a consistent global timing shift.

\paragraph{Duration, pitch, and velocity strata.}
Offset failures below use the direct rule
$|\Delta_{\mathrm{offset}}| >
\max(0.2\,d_{\mathrm{ref}}, 50\text{ ms})$ on the onset+pitch matched pairs.
Intervals are descriptive and use a piece-cluster bootstrap (20,000
resamples of the 137 pieces), so notes from the same piece are not treated
as independent bootstrap units.

\begin{center}
\begin{tabular}{lrrr}
\toprule
\textbf{Reference duration} & \textbf{Pairs} & \textbf{Failure rate} &
\textbf{Piece-bootstrap 95\% interval} \\
\midrule
Under 250 ms & 12,809 & 11.69\% & [9.94\%, 13.69\%] \\
250--500 ms & 1,692 & 43.14\% & [38.11\%, 48.21\%] \\
500 ms--1 s & 825 & 42.18\% & [36.43\%, 48.07\%] \\
At least 1 s & 450 & 41.11\% & [33.55\%, 48.80\%] \\
\bottomrule
\end{tabular}
\end{center}

The offset failure rate also varied with reference velocity, whereas the
three broad pitch groups were more similar:

\begin{center}
\begin{tabular}{lrrr}
\toprule
\textbf{Reference attribute bin} & \textbf{Pairs} & \textbf{Failure rate} &
\textbf{Piece-bootstrap 95\% interval} \\
\midrule
Velocity $\leq 40$ & 1,105 & 37.10\% & [32.92\%, 40.99\%] \\
Velocity 41--80 & 11,033 & 15.59\% & [13.15\%, 18.33\%] \\
Velocity $\geq 81$ & 3,638 & 17.32\% & [14.06\%, 20.85\%] \\
Pitch $\leq 48$ & 1,848 & 15.91\% & [12.75\%, 19.40\%] \\
Pitch 49--72 & 9,446 & 18.15\% & [15.50\%, 21.04\%] \\
Pitch $\geq 73$ & 4,482 & 16.78\% & [13.84\%, 20.26\%] \\
\bottomrule
\end{tabular}
\end{center}

The low-velocity association is still present within duration strata: for
notes shorter than 250 ms, the failure rate was 27.93\% for velocity
$\leq 40$, 10.32\% for velocities 41--80, and 12.52\% for velocities
$\geq 81$. This suggests the low-velocity association is not explained only
by longer durations in this sample. These univariate and stratified
associations do not establish causality and may reflect annotation,
performance, or model effects. The note-level records are retained locally
at \path{checkpoint/maestro-val-offset-note-analysis.json}; the sampled
window manifest and aggregate residuals are at
\path{checkpoint/maestro-val-offset-analysis.json}.

\section{Matched MAESTRO validation comparison}

The original packaged checkpoint and the 10,000-step fine-tuned model were
transcribed on the same deterministic sample windows from all 137 validation
pieces. Each window is 16 seconds; only its saved central 12-second region is
scored. The baseline uses the original checkpoint-era model implementation
(the parent of commit \texttt{354b4c4d}) and all 184 packaged checkpoint
tensors loaded strictly. Both models use the same stereo 44.1 kHz audio and
decoding procedure. Ground truth uses the no-pedal-extension validation
metadata. Metrics are macro-averaged per piece; uncertainty is a paired
bootstrap over pieces (10,000 replicates, seed 20261109). Onset tolerance is
50 ms, and offset matching uses ratio 0.2 with a 50 ms minimum tolerance.

\begin{center}
\begin{tabular}{p{0.25\linewidth}rrp{0.35\linewidth}}
\toprule
\textbf{Metric} & \textbf{Original F1} & \textbf{10k F1} &
\textbf{Paired $\Delta$ F1 (95\% interval)} \\
\midrule
Note onset & 0.964880 & 0.966170 & +0.001290 [-0.000906, +0.003311] \\
Note + offset & 0.753832 & 0.754235 & +0.000403 [-0.006670, +0.007739] \\
Note + velocity & 0.962070 & 0.962087 & +0.000017 [-0.002429, +0.002282] \\
Note + velocity + offset & 0.751670 & 0.751068 & -0.000602 [-0.007605, +0.006691] \\
\bottomrule
\end{tabular}
\end{center}

Supplementary precision and recall summaries were computed from the same
137 per-piece records for the primary note+offset criterion and the
note-onset guard. Paired piece-bootstrap intervals use 10,000 replicates
and seed 20261113; this decomposition was not a separately pre-registered
endpoint.

\begin{center}
\begin{tabular}{llrrl}
\toprule
\textbf{Metric} & \textbf{Quantity} & \textbf{Original} & \textbf{10k} &
\textbf{Paired change (95\% interval)}\\
\midrule
Note onset & Precision & 0.980839 & 0.979504 & -0.001335 [-0.005979, +0.001897]\\
Note onset & Recall & 0.951742 & 0.954816 & +0.003074 [+0.001055, +0.005148]\\
Note+offset & Precision & 0.767418 & 0.765708 & -0.001710 [-0.009022, +0.006017]\\
Note+offset & Recall & 0.742169 & 0.744178 & +0.002009 [-0.005062, +0.009643]\\
\bottomrule
\end{tabular}
\end{center}

These are macro-averaged per-piece precision and recall values. They are
not inputs from which to reconstruct macro F1: F1 was computed for each
piece and then averaged separately.

Across the 137 fixed windows, the 10k model improved, declined, or tied on
58/37/42 pieces for onset, 63/63/11 for note+offset, 55/43/39 for
note+velocity, and 61/68/8 for the combined criterion. The median per-piece
F1 change was zero for all four metrics. For note+offset, the 25th and
75th percentiles were -0.025974 and +0.015464, with a minimum of -0.109040
and maximum of +0.218769. These are descriptive distribution summaries,
not an additional significance test; the largest shifts come from
individual sampled windows and should not be generalized to whole pieces.

The paired note+offset effect is small, and its interval includes both
improvement and degradation. The planning discussion proposed (but did not
pre-register or formally agree on) a practical success threshold of at least
0.005 absolute note+offset F1 gain, a paired interval above zero, and no more
than 0.002 loss in note-onset F1. This comparison does not meet that proposed
threshold. The result therefore supports \textbf{no demonstrated gain};
because only one sampled window per piece was evaluated and the uncertainty
is substantial, it does not establish that the models are equivalent or
that structured attributes cannot help. Do not make a final success claim
from these windows. The complete per-piece results are retained at
\path{checkpoint/maestro-baseline-vs-10000-validation.json}.

\section{Decision and next steps}

1. Do not extend broad fine-tuning or tune on SMD: current validation evidence
   does not establish a practically useful gain.
2. Run one bounded, inference-only offset-score ablation on the same 137
   validation windows. The hypothesis is that moderately increasing the
   offset-timing and offset-presence evidence during interval decoding can
   improve note endpoints without changing note-onset decisions substantially.
   Keep model weights fixed; compare the default factor scale of 1.0 against
   one setting, 1.5. Do not sweep additional values.
3. Before evaluating that setting, define success as at least 0.005 absolute
   note+offset macro-F1 gain over scale 1.0, a paired piece-bootstrap 95\%
   interval wholly above zero, and a note-onset lower confidence bound above
   -0.002. If any condition fails, stop this line without another validation
   tweak; the test set remains reserved.

\subsection{Offset-score ablation result}

The offset-only factor scale was increased from 1.0 to 1.5 during interval
decoding. Model weights and attribute logits were unchanged. This setting
was evaluated on the same 137 windows and compared piece-by-piece with the
scale-1.0 results above; uncertainty uses 10,000 paired piece-bootstrap
replicates with seed 20261110.

\begin{center}
\begin{tabular}{p{0.25\linewidth}rrp{0.35\linewidth}}
\toprule
\textbf{Metric} & \textbf{Scale 1.0 F1} & \textbf{Scale 1.5 F1} &
\textbf{Paired $\Delta$ F1 (95\% interval)} \\
\midrule
Note onset & 0.966170 & 0.966130 & -0.000040 [-0.000650, +0.000624] \\
Note + offset & 0.754235 & 0.754085 & -0.000150 [-0.002370, +0.001996] \\
Note + velocity & 0.962087 & 0.961987 & -0.000100 [-0.000734, +0.000540] \\
Note + velocity + offset & 0.751068 & 0.750956 & -0.000112 [-0.002279, +0.002019] \\
\bottomrule
\end{tabular}
\end{center}

The note+offset result did not meet the pre-set gain or confidence-interval
criteria. The note-onset non-inferiority guard was satisfied, but that alone
is insufficient for success. This targeted inference-only ablation therefore
shows no practical benefit; stop further offset-score tuning and do not use
the MAESTRO test set to rescue or select this setting. The per-piece output
is retained at
\path{checkpoint/maestro-offset-factor-ablation-validation.json}.

The current research result is \textbf{no demonstrated improvement} over the
original model, not proof of equivalence or a universal failure of structured
attributes. Further work should require a new, independently justified
hypothesis and a predeclared criterion before training; otherwise, stop here.

\section{Full-recording validation robustness run}

The fixed-window comparison above scores one 16-second excerpt from each
validation recording. To assess whether the result is sensitive to that
sampling choice, a separate paired evaluation is running over each complete
audio recording. This is a broader coverage of the same validation split,
not a new independent test: the split has already been examined in the
training and sampled-window analyses.

\begin{itemize}
  \item \textbf{Systems:} the original packaged checkpoint, evaluated with
        the historical model implementation from commit
        \texttt{03b23776bae3c8b3046dcbe1f2b17bea64c5b25f}, and the 10k
        structured-attribute checkpoint. Both checkpoints are loaded
        strictly; the evaluator records their SHA-256 hashes.
  \item \textbf{Coverage and segmentation:} all 137 complete MAESTRO
        validation recordings; model inference uses 16-second segments
        with an 8-second hop and retains predictions across the full
        recording.
  \item \textbf{Matching:} no ground-truth pedal extension, 50 ms onset
        tolerance, and offset tolerance equal to the larger of 20\% of
        reference duration or 50 ms.
  \item \textbf{Aggregation and uncertainty:} per-piece F1 is computed for
        each system and then macro-averaged. Paired uncertainty resamples
        pieces (10,000 percentile-bootstrap replicates, seed 20261112).
  \item \textbf{Progress:} the first run was interrupted after 47 of 137
        pieces had been saved. It was resumed from the incremental JSON
        using the same metadata, configuration, checkpoints, device, and
        scoring protocol. The first resume attempt exited before evaluating
        another piece because strict loading reported a candidate
        checkpoint/model-key mismatch. The cause was confirmed: launching
        the file from \texttt{scripts/} placed that folder first on
        \texttt{sys.path}, so Python imported the installed legacy
        \texttt{transkun} package instead of the workspace implementation.
        The evaluator now inserts the repository root before importing
        \texttt{transkun}; the workspace model and candidate checkpoint
        load strictly together. The saved JSON still reports
        \texttt{complete=false}. No full-recording estimate or
        success/failure conclusion is reported until all pieces are
        complete and the output is validated.
\end{itemize}

Because this validation split has informed prior comparisons, the completed
full-recording result should be described as a robustness or sensitivity
analysis, not a confirmatory test of the research hypothesis. Do not adjust
weights, select checkpoints, or launch follow-up training based on its
partial output. The MAESTRO test split remains untouched.

\section{Evidence-based pathways for further research}

The present results do not justify another unstructured parameter sweep.
The following paths answer distinct questions and should be treated as
separate experiments rather than combined into a single claim.

\begin{enumerate}
  \item \textbf{Complete the current robustness analysis.} Validate the
        final JSON for 137 unique records, finite metric values, matching
        protocol and checkpoint hashes, and a completed bootstrap summary.
        Compare it with the fixed-window estimate, but do not use the
        validation result to tune the existing 10k model. If the effects
        differ, report the discrepancy and inspect piece-level changes as
        descriptive diagnostics.
  \item \textbf{Measure whether the new heads learn useful attributes.}
        On a held-out evaluation protocol, report categorical loss,
        calibration, and accuracy for velocity, onset/offset timing, and
        boundary presence separately. Relate those diagnostics to event
        selection errors. This tests the mechanism directly; note F1 alone
        cannot establish whether a head is accurate or calibrated.
  \item \textbf{Separate the attribute contributions.} If the mechanism
        diagnostics support further work, compare a frozen baseline with
        pre-specified velocity-only, timing-only, boundary-presence-only,
        and joint variants. Use matched training budgets and data order
        where practical, and repeat each condition with multiple random
        seeds. Report all runs and uncertainty rather than choosing the
        best seed.
  \item \textbf{Test generalization on independent recordings.} Before any
        final claim, choose an external dataset and define compatible
        event, pedal, and split conventions. Keep recording or performer
        overlap out of train and evaluation partitions. If MAESTRO test is
        used, freeze the selected system and protocol first, then evaluate
        it once as the final held-out result.
  \item \textbf{Stop if the effect is not practically useful.} State an
        absolute minimum worthwhile gain, uncertainty rule, and allowable
        note-onset degradation before the next experiment. A result that
        misses the rule should close that specific configuration family;
        it should not trigger another search against the same validation
        examples.
\end{enumerate}

The current stopping point is therefore bounded: finish and report the
already-started full-recording comparison, then decide whether the evidence
supports a new, independently evaluable experiment. No new training run has
been started.

\section{Artifact locations}

Generated checkpoints, MIDI predictions, and JSON metrics are local and
ignored by Git under \texttt{checkpoint/}. Important result files:

\begin{itemize}
  \item Training-validation TensorBoard event logs:
        \begin{itemize}
          \item Pilot:
                \texttt{checkpoint/structured-attributes-}\\
                \texttt{pilot.pt.log/}
          \item 5k:
                \texttt{checkpoint/structured-attributes-}\\
                \texttt{5000.pt.log/}
          \item 10k:
                \texttt{checkpoint/structured-attributes-}\\
                \texttt{10000.pt.log/}
        \end{itemize}
  \item 1k SMD metrics:
        \path{checkpoint/structured-attributes-pilot-smd/metrics.json}
  \item 5k SMD metrics:
        \path{checkpoint/structured-attributes-5000-smd/metrics.json}
  \item 10k SMD metrics:
        \path{checkpoint/structured-attributes-10000-smd/metrics.json}
  \item Matched original-versus-10k MAESTRO validation comparison:
        \\
        \texttt{checkpoint/maestro-baseline-}\\
        \texttt{vs-10000-validation.json}
  \item Incremental full-recording original-versus-10k comparison:
        \\
        \texttt{checkpoint/maestro-baseline-}\\
        \texttt{vs-10000-full-validation.json}
  \item Full-recording evaluation driver:
        \path{scripts/evaluate_maestro_full_validation.py}
  \item Targeted offset-score ablation:
        \\
        \texttt{checkpoint/maestro-offset-factor-}\\
        \texttt{ablation-validation.json}
  \item Weight 0 SMD metrics:
        \\
        \texttt{checkpoint/structured-attributes-10000-}\\
        \texttt{no-attribute-score-smd/metrics.json}
  \item Weight 0.5 SMD metrics:
        \path{checkpoint/structured-attributes-10000-weight-0.5-smd/metrics.json}
  \item Weight 0.75 SMD metrics:
        \\
        \texttt{checkpoint/structured-attributes-10000-}\\
        \texttt{weight-0.75-smd/metrics.json}
\end{itemize}

\section{References}

\begin{itemize}
  \item User-provided OpenReview paper:
        \url{https://openreview.net/pdf?id=DGA8XbJ8FVd}
  \item User-provided arXiv paper:
        \url{https://arxiv.org/pdf/2404.09466}
  \item Repository and project-specific evaluation implementation:
        \texttt{transkun/Evaluation.py} and
        \texttt{transkun/computeMetrics.py}.
\end{itemize}

\end{document}
